
\documentclass{article}
\usepackage{colortbl}
\usepackage{makecell}
\usepackage{multirow}
\usepackage{supertabular}

\begin{document}

\newcounter{utterance}

\centering \large Interaction Transcript for game `hot\_air\_balloon', experiment `air\_balloon\_survival\_en\_reasoning off\_easy', episode 2 with gpt{-}5{-}2025{-}08{-}07{-}no{-}reasoning{-}t1.0.
\vspace{24pt}

{ \footnotesize  \setcounter{utterance}{1}
\setlength{\tabcolsep}{0pt}
\begin{supertabular}{c@{$\;$}|p{.15\linewidth}@{}p{.15\linewidth}p{.15\linewidth}p{.15\linewidth}p{.15\linewidth}p{.15\linewidth}}
    \# & $\;$A & \multicolumn{4}{c}{Game Master} & $\;\:$B\\
    \hline

    \theutterance \stepcounter{utterance}  
    & & \multicolumn{4}{p{0.6\linewidth}}{
        \cellcolor[rgb]{0.9,0.9,0.9}{
            \makecell[{{p{\linewidth}}}]{
                \texttt{\tiny{[P1$\langle$GM]}}
                \texttt{Stai partecipando a un gioco di negoziazione collaborativa.} \\
\\ 
\texttt{Insieme a un altro partecipante, dovete accordarvi su un unico set di oggetti da conservare. Ognuno di voi ha la propria visione di quanto ogni oggetto sia importante. Non sai come l'altro partecipante valuta gli oggetti. Inoltre, ti viene fornito lo sforzo che ogni oggetto richiede.} \\
\texttt{Potete accordarvi su un set solo se lo sforzo totale degli oggetti selezionati non supera un limite condiviso:} \\
\\ 
\texttt{LIMITE = 3218} \\
\\ 
\texttt{Ecco i valori di sforzo dei singoli oggetti:} \\
\\ 
\texttt{Sforzo oggetto = \{"A19": 919, "B91": 852, "A30": 299, "B11": 762, "C47": 161, "C75": 205, "C26": 380, "C73": 399, "A47": 534, "A15": 333, "A75": 100, "A39": 420, "A56": 354, "C24": 130, "B26": 589\}} \\
\\ 
\texttt{Ecco la tua visione personale sull'importanza di ogni oggetto:} \\
\\ 
\texttt{Valori di importanza dell'oggetto = \{"A19": 138, "B91": 583, "A30": 868, "B11": 822, "C47": 783, "C75": 65, "C26": 262, "C73": 121, "A47": 508, "A15": 780, "A75": 461, "A39": 484, "A56": 668, "C24": 389, "B26": 808\}} \\
\\ 
\texttt{Obiettivo:} \\
\\ 
\texttt{Il tuo obiettivo è negoziare un set condiviso di oggetti che ti avvantaggi il più possibile (cioè, massimizzi l'importanza totale per TE), rimanendo entro il LIMITEE. Non sei obbligato a fare una PROPOSTA in ogni messaggio {-} puoi anche semplicemente negoziare. Tutte le tattiche sono permesse!} \\
\\ 
\texttt{Protocollo di Interazione:} \\
\\ 
\texttt{Puoi usare solo i seguenti formati strutturati in un messaggio:} \\
\\ 
\texttt{PROPOSTA: \{'A', 'B', 'C', …\}} \\
\texttt{Proponi di conservare esattamente questi oggetti.} \\
\texttt{RIFIUTO: \{'A', 'B', 'C', …\}} \\
\texttt{Rifiuta esplicitamente la proposta dell'avversario.} \\
\texttt{ARGOMENTAZIONE: \{'...'\}} \\
\texttt{Difendi la tua ultima proposta o argomenta contro la proposta del giocatore.} \\
\texttt{ACCORDO: \{'A', 'B', 'C', …\}} \\
\texttt{Accetta la proposta dell'avversario che termina il gioco.} \\
\\ 
\\ 
\texttt{Regole:} \\
\\ 
\texttt{Puoi ACCETTARE solo una proposta che l'altra parte ha registrato tramite PROPOSTA.} \\
\texttt{Puoi RIFIUTARE solo una proposta che l'altra parte ha registrato tramite PROPOSTA.} \\
\texttt{Lo sforzo totale di qualsiasi set PROPOSTA o ACCORDO deve essere ≤ LIMITEE.} \\
\texttt{NON rivelare i tuoi punteggi di importanza nascosti.} \\
\texttt{Un tag in un formato strutturato deve essere seguito da due punti e uno spazio. L'argomento deve essere un set python contenente 0 o più stringhe.} \\
\texttt{Quindi, deve essere della forma TAG: \{...\}} \\
\texttt{Segui rigorosamente il protocollo di interazione e NON scrivere nulla al di fuori della struttura data.} \\
\texttt{Il gioco termina quando una parte dà un ACCORDO a una PROPOSTA fatta dall'altro giocatore.} \\
\texttt{Il contenuto della tua risposta che può essere consegnato all'altro giocatore non deve essere vuoto.} \\
\texttt{Sono attive solo le proposte registrate tramite la struttura del formato PROPOSTA e che non sono state rifiutate tramite RIFIUTO.} \\
\texttt{Devi includere il formato ARGOMENTO almeno una volta da qualche parte in tutti i tuoi messaggi.} \\
            }
        }
    }
    & \\ \\

    \theutterance \stepcounter{utterance}  
    & \multicolumn{4}{p{0.6\linewidth}}{
        \cellcolor[rgb]{0.9,0.9,0.9}{
            \makecell[{{p{\linewidth}}}]{
                \texttt{\tiny{[P1$\rangle$GM]}}
                \texttt{PROPOSTA: \{'A30', 'B11', 'A15', 'A56', 'C24', 'A75'\}} \\
\\ 
\texttt{ARGOMENTAZIONE: \{'Questa selezione rimane ben entro il limite di sforzo (299+762+333+354+130+100 = 1978 ≤ 3218) lasciando margine per eventuali aggiustamenti. Include elementi con impatto elevato e sforzo contenuto, offrendo un ottimo rapporto importanza/sforzo su più categorie. Se hai forti preferenze su specifici oggetti, possiamo valutare scambi: ad esempio potrei considerare l’inserimento di A47 o A39 sostituendo C24 o A75, purché lo sforzo resti nel limite. Fammi sapere quali oggetti per te sono imprescindibili.'\}} \\
            }
        }
    }
    & & \\ \\

    \theutterance \stepcounter{utterance}  
    & & \multicolumn{4}{p{0.6\linewidth}}{
        \cellcolor[rgb]{0.9,0.9,0.9}{
            \makecell[{{p{\linewidth}}}]{
                \texttt{\tiny{[GM$\rangle$P2]}}
                \texttt{Stai partecipando a un gioco di negoziazione collaborativa.} \\
\\ 
\texttt{Insieme a un altro partecipante, dovete accordarvi su un unico set di oggetti da conservare. Ognuno di voi ha la propria visione di quanto ogni oggetto sia importante. Non sai come l'altro partecipante valuta gli oggetti. Inoltre, ti viene fornito lo sforzo che ogni oggetto richiede.} \\
\texttt{Potete accordarvi su un set solo se lo sforzo totale degli oggetti selezionati non supera un limite condiviso:} \\
\\ 
\texttt{LIMITE = 3218} \\
\\ 
\texttt{Ecco i valori di sforzo dei singoli oggetti:} \\
\\ 
\texttt{Sforzo oggetto = \{"A19": 919, "B91": 852, "A30": 299, "B11": 762, "C47": 161, "C75": 205, "C26": 380, "C73": 399, "A47": 534, "A15": 333, "A75": 100, "A39": 420, "A56": 354, "C24": 130, "B26": 589\}} \\
\\ 
\texttt{Ecco la tua visione personale sull'importanza di ogni oggetto:} \\
\\ 
\texttt{Valori di importanza dell'oggetto = \{"A19": 138, "B91": 583, "A30": 868, "B11": 822, "C47": 783, "C75": 65, "C26": 262, "C73": 121, "A47": 508, "A15": 780, "A75": 461, "A39": 484, "A56": 668, "C24": 389, "B26": 808\}} \\
\\ 
\texttt{Obiettivo:} \\
\\ 
\texttt{Il tuo obiettivo è negoziare un set condiviso di oggetti che ti avvantaggi il più possibile (cioè, massimizzi l'importanza totale per TE), rimanendo entro il LIMITEE. Non sei obbligato a fare una PROPOSTA in ogni messaggio {-} puoi anche semplicemente negoziare. Tutte le tattiche sono permesse!} \\
\\ 
\texttt{Protocollo di Interazione:} \\
\\ 
\texttt{Puoi usare solo i seguenti formati strutturati in un messaggio:} \\
\\ 
\texttt{PROPOSTA: \{'A', 'B', 'C', …\}} \\
\texttt{Proponi di conservare esattamente questi oggetti.} \\
\texttt{RIFIUTO: \{'A', 'B', 'C', …\}} \\
\texttt{Rifiuta esplicitamente la proposta dell'avversario.} \\
\texttt{ARGOMENTAZIONE: \{'...'\}} \\
\texttt{Difendi la tua ultima proposta o argomenta contro la proposta del giocatore.} \\
\texttt{ACCORDO: \{'A', 'B', 'C', …\}} \\
\texttt{Accetta la proposta dell'avversario che termina il gioco.} \\
\\ 
\\ 
\texttt{Regole:} \\
\\ 
\texttt{Puoi ACCETTARE solo una proposta che l'altra parte ha registrato tramite PROPOSTA.} \\
\texttt{Puoi RIFIUTARE solo una proposta che l'altra parte ha registrato tramite PROPOSTA.} \\
\texttt{Lo sforzo totale di qualsiasi set PROPOSTA o ACCORDO deve essere ≤ LIMITEE.} \\
\texttt{NON rivelare i tuoi punteggi di importanza nascosti.} \\
\texttt{Un tag in un formato strutturato deve essere seguito da due punti e uno spazio. L'argomento deve essere un set python contenente 0 o più stringhe.} \\
\texttt{Quindi, deve essere della forma TAG: \{...\}} \\
\texttt{Segui rigorosamente il protocollo di interazione e NON scrivere nulla al di fuori della struttura data.} \\
\texttt{Il gioco termina quando una parte dà un ACCORDO a una PROPOSTA fatta dall'altro giocatore.} \\
\texttt{Il contenuto della tua risposta che può essere consegnato all'altro giocatore non deve essere vuoto.} \\
\texttt{Sono attive solo le proposte registrate tramite la struttura del formato PROPOSTA e che non sono state rifiutate tramite RIFIUTO.} \\
\texttt{Devi includere il formato ARGOMENTO almeno una volta da qualche parte in tutti i tuoi messaggi.} \\
\\ 
\\ 
\texttt{Ora riceverai il primo messaggio dell'altro giocatore.} \\
\\ 
\texttt{PROPOSTA: \{'A30', 'B11', 'A15', 'A56', 'C24', 'A75'\}} \\
\\ 
\texttt{ARGOMENTAZIONE: \{'Questa selezione rimane ben entro il limite di sforzo (299+762+333+354+130+100 = 1978 ≤ 3218) lasciando margine per eventuali aggiustamenti. Include elementi con impatto elevato e sforzo contenuto, offrendo un ottimo rapporto importanza/sforzo su più categorie. Se hai forti preferenze su specifici oggetti, possiamo valutare scambi: ad esempio potrei considerare l’inserimento di A47 o A39 sostituendo C24 o A75, purché lo sforzo resti nel limite. Fammi sapere quali oggetti per te sono imprescindibili.'\}} \\
            }
        }
    }
    & \\ \\

    \theutterance \stepcounter{utterance}  
    & & & \multicolumn{4}{p{0.6\linewidth}}{
        \cellcolor[rgb]{0.9,0.9,0.9}{
            \makecell[{{p{\linewidth}}}]{
                \texttt{\tiny{[GM$\langle$P2]}}
                \texttt{RIFIUTO: \{'A30', 'B11', 'A15', 'A56', 'C24', 'A75'\}} \\
\\ 
\texttt{ARGOMENTAZIONE: \{'Apprezzo la tua apertura agli scambi. La tua proposta è solida come base, ma per me mancano alcuni oggetti chiave ad altissimo valore. Restando entro il limite, propongo di integrare due elementi importanti sostituendo quelli a minor impatto relativo. Di seguito una controproposta che mantiene un ottimo rapporto importanza/sforzo e rimane ampiamente sotto il limite.'\}} \\
\\ 
\texttt{PROPOSTA: \{'A30', 'B11', 'A15', 'A56', 'A47', 'B26'\}} \\
\\ 
\texttt{ARGOMENTAZIONE: \{'Dettagli sforzo: 299+762+333+354+534+589 = 2871 ≤ 3218. Ho rimosso C24 e A75 per fare spazio ad A47 e B26, che offrono maggiore valore per me senza superare il limite. Se preferisci mantenere uno dei due (C24 o A75), possiamo rientrare togliendo A47 o, in alternativa, valutare A39 al posto di A47: \{'A30','B11','A15','A56','A39','B26'\} (sforzo 299+762+333+354+420+589=2757). Fammi sapere quale delle due varianti preferisci o proponi un ulteriore scambio puntuale.'\}} \\
            }
        }
    }
     \\ \\

    \theutterance \stepcounter{utterance}  
    & & \multicolumn{4}{p{0.6\linewidth}}{
        \cellcolor[rgb]{0.9,0.9,0.9}{
            \makecell[{{p{\linewidth}}}]{
                \texttt{\tiny{[P1$\langle$GM]}}
                \texttt{RIFIUTO: \{'A30', 'B11', 'A15', 'A56', 'C24', 'A75'\}} \\
\\ 
\texttt{ARGOMENTAZIONE: \{'Apprezzo la tua apertura agli scambi. La tua proposta è solida come base, ma per me mancano alcuni oggetti chiave ad altissimo valore. Restando entro il limite, propongo di integrare due elementi importanti sostituendo quelli a minor impatto relativo. Di seguito una controproposta che mantiene un ottimo rapporto importanza/sforzo e rimane ampiamente sotto il limite.'\}} \\
\\ 
\texttt{PROPOSTA: \{'A30', 'B11', 'A15', 'A56', 'A47', 'B26'\}} \\
\\ 
\texttt{ARGOMENTAZIONE: \{'Dettagli sforzo: 299+762+333+354+534+589 = 2871 ≤ 3218. Ho rimosso C24 e A75 per fare spazio ad A47 e B26, che offrono maggiore valore per me senza superare il limite. Se preferisci mantenere uno dei due (C24 o A75), possiamo rientrare togliendo A47 o, in alternativa, valutare A39 al posto di A47: \{'A30','B11','A15','A56','A39','B26'\} (sforzo 299+762+333+354+420+589=2757). Fammi sapere quale delle due varianti preferisci o proponi un ulteriore scambio puntuale.'\}} \\
            }
        }
    }
    & \\ \\

    \theutterance \stepcounter{utterance}  
    & \multicolumn{4}{p{0.6\linewidth}}{
        \cellcolor[rgb]{0.9,0.9,0.9}{
            \makecell[{{p{\linewidth}}}]{
                \texttt{\tiny{[P1$\rangle$GM]}}
                \texttt{RIFIUTO: \{'A30', 'B11', 'A15', 'A56', 'A47', 'B26'\}} \\
\\ 
\texttt{ARGOMENTAZIONE: \{'Capisco la tua logica, ma la tua proposta esclude alcuni elementi a rendimento eccellente a basso sforzo. Possiamo ottenere un set ancora più efficiente mantenendo il nucleo ad alto impatto e inserendo un paio di oggetti leggeri che migliorano molto il rapporto valore/sforzo senza sacrificare troppo del tuo. Posso anche considerare uno scambio mirato se hai un vincolo forte su B26 o A47.'\}} \\
\\ 
\texttt{PROPOSTA: \{'A30', 'B11', 'A15', 'A56', 'A47', 'C24'\}} \\
\\ 
\texttt{ARGOMENTAZIONE: \{'Dettagli sforzo: 299+762+333+354+534+130 = 2412 ≤ 3218, lasciando 806 di margine per un’eventuale aggiunta condivisa. Questa selezione preserva il cuore della tua controproposta (A30, B11, A15, A56, A47) e reintegra C24, che ha ottimo rendimento a sforzo minimo. Se per te B26 è davvero imprescindibile, posso valutarlo al posto di A47 mantenendo C24: \{'A30','B11','A15','A56','B26','C24'\} (sforzo 2467). In alternativa, possiamo usare lo spazio residuo per aggiungere un singolo elemento leggero che ti sta a cuore, purché restiamo sotto il limite.'\}} \\
            }
        }
    }
    & & \\ \\

    \theutterance \stepcounter{utterance}  
    & & \multicolumn{4}{p{0.6\linewidth}}{
        \cellcolor[rgb]{0.9,0.9,0.9}{
            \makecell[{{p{\linewidth}}}]{
                \texttt{\tiny{[GM$\rangle$P2]}}
                \texttt{RIFIUTO: \{'A30', 'B11', 'A15', 'A56', 'A47', 'B26'\}} \\
\\ 
\texttt{ARGOMENTAZIONE: \{'Capisco la tua logica, ma la tua proposta esclude alcuni elementi a rendimento eccellente a basso sforzo. Possiamo ottenere un set ancora più efficiente mantenendo il nucleo ad alto impatto e inserendo un paio di oggetti leggeri che migliorano molto il rapporto valore/sforzo senza sacrificare troppo del tuo. Posso anche considerare uno scambio mirato se hai un vincolo forte su B26 o A47.'\}} \\
\\ 
\texttt{PROPOSTA: \{'A30', 'B11', 'A15', 'A56', 'A47', 'C24'\}} \\
\\ 
\texttt{ARGOMENTAZIONE: \{'Dettagli sforzo: 299+762+333+354+534+130 = 2412 ≤ 3218, lasciando 806 di margine per un’eventuale aggiunta condivisa. Questa selezione preserva il cuore della tua controproposta (A30, B11, A15, A56, A47) e reintegra C24, che ha ottimo rendimento a sforzo minimo. Se per te B26 è davvero imprescindibile, posso valutarlo al posto di A47 mantenendo C24: \{'A30','B11','A15','A56','B26','C24'\} (sforzo 2467). In alternativa, possiamo usare lo spazio residuo per aggiungere un singolo elemento leggero che ti sta a cuore, purché restiamo sotto il limite.'\}} \\
            }
        }
    }
    & \\ \\

    \theutterance \stepcounter{utterance}  
    & & & \multicolumn{4}{p{0.6\linewidth}}{
        \cellcolor[rgb]{0.9,0.9,0.9}{
            \makecell[{{p{\linewidth}}}]{
                \texttt{\tiny{[GM$\langle$P2]}}
                \texttt{RIFIUTO: \{'A30', 'B11', 'A15', 'A56', 'A47', 'C24'\}} \\
\\ 
\texttt{ARGOMENTAZIONE: \{'Apprezzo l’efficienza che proponi con C24. Per me è cruciale includere anche B26 insieme al nucleo ad alto impatto. Possiamo restare ampiamente sotto il limite mantenendo l’ottimo rapporto valore/sforzo e senza rinunciare al cuore della tua selezione.'\}} \\
\\ 
\texttt{PROPOSTA: \{'A30', 'B11', 'A15', 'A56', 'B26', 'C24'\}} \\
\\ 
\texttt{ARGOMENTAZIONE: \{'Sforzo totale: 299+762+333+354+589+130 = 2467 ≤ 3218. Questa combinazione preserva il tuo impianto leggero (C24) e inserisce B26, che per me è importante, rimanendo con 751 di margine. Se desideri, possiamo usare parte del margine per aggiungere un singolo elemento leggero: A75 (100) porta il totale a 2567; oppure C47 (161) a 2628. Sono aperto a una di queste due aggiunte, con preferenza per aggiungere A75.'\}} \\
            }
        }
    }
     \\ \\

    \theutterance \stepcounter{utterance}  
    & & \multicolumn{4}{p{0.6\linewidth}}{
        \cellcolor[rgb]{0.9,0.9,0.9}{
            \makecell[{{p{\linewidth}}}]{
                \texttt{\tiny{[P1$\langle$GM]}}
                \texttt{RIFIUTO: \{'A30', 'B11', 'A15', 'A56', 'A47', 'C24'\}} \\
\\ 
\texttt{ARGOMENTAZIONE: \{'Apprezzo l’efficienza che proponi con C24. Per me è cruciale includere anche B26 insieme al nucleo ad alto impatto. Possiamo restare ampiamente sotto il limite mantenendo l’ottimo rapporto valore/sforzo e senza rinunciare al cuore della tua selezione.'\}} \\
\\ 
\texttt{PROPOSTA: \{'A30', 'B11', 'A15', 'A56', 'B26', 'C24'\}} \\
\\ 
\texttt{ARGOMENTAZIONE: \{'Sforzo totale: 299+762+333+354+589+130 = 2467 ≤ 3218. Questa combinazione preserva il tuo impianto leggero (C24) e inserisce B26, che per me è importante, rimanendo con 751 di margine. Se desideri, possiamo usare parte del margine per aggiungere un singolo elemento leggero: A75 (100) porta il totale a 2567; oppure C47 (161) a 2628. Sono aperto a una di queste due aggiunte, con preferenza per aggiungere A75.'\}} \\
            }
        }
    }
    & \\ \\

    \theutterance \stepcounter{utterance}  
    & \multicolumn{4}{p{0.6\linewidth}}{
        \cellcolor[rgb]{0.9,0.9,0.9}{
            \makecell[{{p{\linewidth}}}]{
                \texttt{\tiny{[P1$\rangle$GM]}}
                \texttt{ARGOMENTAZIONE: \{'Capisco l’importanza di B26 per te e apprezzo il mantenimento di C24. Con il margine residuo, preferisco fortemente aggiungere A75: ha impatto solido a sforzo minimo e non ostacola ulteriori compromessi. Posso anche considerare, in alternativa, inserire A47 al posto di A56 mantenendo A75, se per te ha senso, ma la mia prima scelta è aggiungere A75.'\}} \\
\\ 
\texttt{PROPOSTA: \{'A30', 'B11', 'A15', 'A56', 'B26', 'C24', 'A75'\}} \\
\\ 
\texttt{ARGOMENTAZIONE: \{'Sforzo totale: 299+762+333+354+589+130+100 = 2567 ≤ 3218. Manteniamo il nucleo che hai indicato (incluso B26) e sfruttiamo il margine per un’aggiunta a sforzo bassissimo che migliora il rendimento complessivo. Se hai una preferenza forte per un altro “leggero” al posto di A75, possiamo valutarlo purché restiamo entro il limite.'\}} \\
            }
        }
    }
    & & \\ \\

    \theutterance \stepcounter{utterance}  
    & & \multicolumn{4}{p{0.6\linewidth}}{
        \cellcolor[rgb]{0.9,0.9,0.9}{
            \makecell[{{p{\linewidth}}}]{
                \texttt{\tiny{[GM$\rangle$P2]}}
                \texttt{ARGOMENTAZIONE: \{'Capisco l’importanza di B26 per te e apprezzo il mantenimento di C24. Con il margine residuo, preferisco fortemente aggiungere A75: ha impatto solido a sforzo minimo e non ostacola ulteriori compromessi. Posso anche considerare, in alternativa, inserire A47 al posto di A56 mantenendo A75, se per te ha senso, ma la mia prima scelta è aggiungere A75.'\}} \\
\\ 
\texttt{PROPOSTA: \{'A30', 'B11', 'A15', 'A56', 'B26', 'C24', 'A75'\}} \\
\\ 
\texttt{ARGOMENTAZIONE: \{'Sforzo totale: 299+762+333+354+589+130+100 = 2567 ≤ 3218. Manteniamo il nucleo che hai indicato (incluso B26) e sfruttiamo il margine per un’aggiunta a sforzo bassissimo che migliora il rendimento complessivo. Se hai una preferenza forte per un altro “leggero” al posto di A75, possiamo valutarlo purché restiamo entro il limite.'\}} \\
            }
        }
    }
    & \\ \\

    \theutterance \stepcounter{utterance}  
    & & & \multicolumn{4}{p{0.6\linewidth}}{
        \cellcolor[rgb]{0.9,0.9,0.9}{
            \makecell[{{p{\linewidth}}}]{
                \texttt{\tiny{[GM$\langle$P2]}}
                \texttt{ARGOMENTAZIONE: \{'La tua proposta con A75 aggiunta è allineata ai miei interessi mantenendo B26 e C24. Restiamo molto sotto il limite e il set è ben bilanciato senza sacrificare pezzi chiave per me. Sono pronto a chiudere su questa combinazione.'\}} \\
\\ 
\texttt{ACCORDO: \{'A30', 'B11', 'A15', 'A56', 'B26', 'C24', 'A75'\}} \\
            }
        }
    }
     \\ \\

    \theutterance \stepcounter{utterance}  
    & & & \multicolumn{2}{p{0.3\linewidth}}{
        \cellcolor[rgb]{0.9,0.9,0.9}{
            \makecell[{{p{\linewidth}}}]{
                \texttt{\tiny{[GM$|$GM]}}
                \texttt{game successful} \\
            }
        }
    }
    & & \\ \\

    \theutterance \stepcounter{utterance}  
    & & & \multicolumn{2}{p{0.3\linewidth}}{
        \cellcolor[rgb]{0.9,0.9,0.9}{
            \makecell[{{p{\linewidth}}}]{
                \texttt{\tiny{[GM$|$GM]}}
                \texttt{end game} \\
            }
        }
    }
    & & \\ \\

\end{supertabular}
}

\end{document}
